\documentclass[11pt]{article}
\usepackage[margin=1in]{geometry}
\usepackage{amsmath}
\usepackage{graphicx}
\usepackage{booktabs}
\usepackage{xcolor}
\usepackage{fancyhdr}
\usepackage{enumitem}

% Define OSU orange color
\definecolor{osaorange}{RGB}{220,115,38}

% Set up headers and footers
\pagestyle{fancy}
\fancyhf{}
\lhead{\textcolor{osaorange}{H524 Introduction to Biostatistics}}
\rhead{Week 10 Worksheet}
\cfoot{\thepage}

\title{\textcolor{osaorange}{\textbf{Week 10 Partner Worksheet}}\\
\large{Regression Diagnostics and Model Assessment}}
\author{}
\date{Fall 2025}

\begin{document}

\maketitle

\noindent\textbf{Names:} \underline{\hspace{6cm}}

\vspace{0.5cm}

\noindent\textbf{Instructions:} Work with your partner to complete these problems. Show your work and explain your reasoning. You may use R for calculations and creating plots.

\section*{Problem 1: Interpreting Diagnostic Plots (25 points)}

A researcher studies the relationship between hours of study per week (X) and exam score (Y) for 80 students. After fitting a simple linear regression, they obtain the following diagnostic plots:

\vspace{0.3cm}

\textbf{Plot A: Residuals vs Fitted}
\begin{itemize}
\item Shows a clear U-shaped pattern
\item Residuals are positive for low fitted values, negative for middle values, positive for high values
\end{itemize}

\textbf{Plot B: Q-Q Plot}
\begin{itemize}
\item Points fall approximately on a straight line
\item Slight deviation at the extreme ends
\end{itemize}

\textbf{Plot C: Scale-Location Plot}
\begin{itemize}
\item Shows relatively horizontal trend line
\item Random scatter of points
\end{itemize}

\textbf{Plot D: Residuals vs Leverage}
\begin{itemize}
\item Three points are labeled: observations 15, 42, and 73
\item Observation 42 has high leverage and falls outside Cook's distance contour
\end{itemize}

\vspace{0.5cm}

\textbf{a)} (6 points) Based on the Residuals vs Fitted plot, which regression assumption is violated? What does this pattern suggest about the relationship between study hours and exam score?

\vspace{3cm}

\textbf{b)} (5 points) What does the Q-Q plot tell you about the normality assumption? Should you be concerned?

\vspace{2.5cm}

\textbf{c)} (5 points) What does the Scale-Location plot indicate about the homoscedasticity assumption?

\vspace{2.5cm}

\textbf{d)} (5 points) Observation 42 appears to be influential. What should the researcher do about this observation? List at least three steps.

\vspace{3cm}

\textbf{e)} (4 points) Given the U-shaped pattern in the residuals, suggest two potential remedies to improve the model fit.

\vspace{2.5cm}

\newpage

\section*{Problem 2: Cook's Distance and Influence (20 points)}

A study examines the relationship between age (X) and bone mineral density (Y) in 50 postmenopausal women. The regression yields:

$$\hat{Y} = 1.25 - 0.015X, \quad R^2 = 0.42$$

Cook's distances for selected observations are:

\begin{center}
\begin{tabular}{cccc}
\toprule
Observation & Age (years) & BMD (g/cm$^2$) & Cook's D \\
\midrule
1 & 52 & 0.98 & 0.02 \\
12 & 68 & 0.23 & 0.78 \\
27 & 59 & 0.82 & 0.05 \\
38 & 75 & 0.12 & 0.12 \\
49 & 54 & 1.15 & 0.03 \\
\bottomrule
\end{tabular}
\end{center}

\vspace{0.3cm}

\textbf{a)} (5 points) Which observation(s) are influential based on Cook's distance? Use the threshold of 0.5.

\vspace{2cm}

\textbf{b)} (6 points) Looking at observation 12, explain why it has high Cook's distance. Consider both its age and BMD values relative to the regression line.

\vspace{3cm}

\textbf{c)} (5 points) The researcher refits the model without observation 12 and obtains:

$$\hat{Y} = 1.18 - 0.012X, \quad R^2 = 0.61$$

Compare this to the original model. How did removing the influential point affect:
\begin{itemize}
\item The slope?
\item The R-squared?
\end{itemize}

\vspace{3cm}

\textbf{d)} (4 points) Should the researcher report only the model without observation 12, only the model with it, or both? Explain your reasoning.

\vspace{3cm}

\newpage

\section*{Problem 3: Transformations (25 points)}

A pharmaceutical company studies the relationship between drug dose (mg) and plasma concentration ($\mu$g/mL) in 60 patients. The data show strong right skewness in plasma concentration. Two models are fit:

\textbf{Model 1 (Original scale):}
$$\widehat{\text{Concentration}} = 2.5 + 1.8 \times \text{Dose}, \quad R^2 = 0.35$$

Diagnostic plots show:
\begin{itemize}
\item Funnel-shaped pattern in residuals (variance increases with fitted values)
\item Right-skewed residuals in Q-Q plot
\end{itemize}

\textbf{Model 2 (Log-transformed):}
$$\widehat{\log(\text{Concentration})} = 0.85 + 0.12 \times \text{Dose}, \quad R^2 = 0.72$$

Diagnostic plots show:
\begin{itemize}
\item Random scatter in residuals
\item Nearly normal residuals in Q-Q plot
\end{itemize}

\vspace{0.5cm}

\textbf{a)} (6 points) Which model better satisfies the regression assumptions? Justify your answer based on the diagnostic information provided.

\vspace{3cm}

\textbf{b)} (7 points) For Model 2, interpret the slope of 0.12 in practical terms. Specifically, if the dose increases by 1 mg, by what percentage does the plasma concentration change?

\textit{Hint: Use the formula $(e^{\beta_1} - 1) \times 100\%$}

\vspace{3cm}

\textbf{c)} (6 points) Using Model 2, predict the plasma concentration for a patient receiving a 10 mg dose. Show your calculation.

\textit{Hint: First predict log(Concentration), then back-transform using $e^{\text{prediction}}$}

\vspace{3cm}

\textbf{d)} (6 points) A colleague suggests trying a square root transformation instead of log. Under what circumstances would a square root transformation be more appropriate than a log transformation?

\vspace{3cm}

\newpage

\section*{Problem 4: Confidence vs Prediction Intervals (15 points)}

A cardiologist studies systolic blood pressure (SBP) as a function of age. Based on 100 patients, the fitted model is:

$$\widehat{\text{SBP}} = 95 + 0.8 \times \text{Age}$$

with residual standard error $s_e = 12$ mmHg.

For a 60-year-old patient:
\begin{itemize}
\item 95\% Confidence Interval for mean SBP: (131.6, 136.4) mmHg
\item 95\% Prediction Interval for individual SBP: (108.2, 159.8) mmHg
\end{itemize}

\vspace{0.5cm}

\textbf{a)} (4 points) What is the predicted SBP for a 60-year-old patient?

\vspace{2cm}

\textbf{b)} (6 points) Explain the difference between the confidence interval and prediction interval in this context. Why is the prediction interval wider?

\vspace{3.5cm}

\textbf{c)} (5 points) A physician wants to know: ``What is the typical blood pressure for 60-year-old patients?'' Which interval should they use? What about if they want to predict the blood pressure for a specific 60-year-old patient coming in tomorrow?

\vspace{3.5cm}

\newpage

\section*{Problem 5: Reporting Results (15 points)}

You are asked to review a colleague's regression analysis. They provide the following summary:

\begin{quote}
\textit{``We found that X is significantly related to Y (p < 0.001). The R-squared is 0.18, which is low but acceptable. The slope is 2.3.''}
\end{quote}

The full model output shows:
\begin{itemize}
\item $n = 250$ observations
\item Slope: $b_1 = 2.3$, $SE = 0.4$
\item 95\% CI for slope: (1.5, 3.1)
\item Intercept: $b_0 = 15.6$
\item $R^2 = 0.18$, $p < 0.001$
\end{itemize}

Diagnostic plots were not shown.

\vspace{0.5cm}

\textbf{a)} (6 points) List three important pieces of information that are missing from the colleague's summary.

\vspace{3cm}

\textbf{b)} (5 points) Rewrite the results in a more complete and informative way, including:
\begin{itemize}
\item Sample size
\item Effect size with confidence interval
\item Interpretation in context (assume Y = cholesterol in mg/dL, X = BMI)
\item Model fit
\end{itemize}

\vspace{4cm}

\textbf{c)} (4 points) Why should the colleague have included diagnostic plots in their report? What could go wrong if assumptions are not checked?

\vspace{3cm}

\newpage

\section*{Bonus Problem: Polynomial Regression (10 points)}

A developmental psychologist studies the relationship between age (3-18 years) and vocabulary score. A scatterplot suggests an inverted U-shape: vocabulary increases rapidly in childhood, peaks around age 12-14, then slightly decreases in adolescence (possibly due to test design).

A quadratic model is fit:

$$\widehat{\text{Score}} = 20 + 15 \times \text{Age} - 0.5 \times \text{Age}^2$$

with $R^2 = 0.78$ and all coefficients significant at $p < 0.01$.

\vspace{0.5cm}

\textbf{a)} (4 points) At what age is vocabulary score predicted to be maximum?

\textit{Hint: The maximum occurs at Age = $-\beta_1 / (2\beta_2)$}

\vspace{3cm}

\textbf{b)} (3 points) Predict the vocabulary score for:
\begin{itemize}
\item A 6-year-old child
\item A 15-year-old adolescent
\end{itemize}

\vspace{3cm}

\textbf{c)} (3 points) The researcher wants to use this model to predict vocabulary scores for 25-year-old adults. What would you advise? Why?

\vspace{3cm}

\section*{Summary Questions}

After completing this worksheet, discuss with your partner:

\begin{enumerate}
\item What is the most important diagnostic plot and why?
\item When should you consider transforming your data?
\item How do you decide whether to remove an influential point?
\item What's the difference between statistical and practical significance?
\end{enumerate}

\vspace{0.5cm}

\hrule

\vspace{0.3cm}

\textbf{Reminder:} For your final project presentations next week:
\begin{itemize}
\item Always check regression assumptions
\item Report confidence intervals, not just p-values
\item Create clear visualizations
\item Interpret results in context
\item Acknowledge limitations
\end{itemize}

\end{document}
